\chapter{Conclusion and Outlook}\label{chap:zusfas}
The aim of this work was to implement a new active learning (AL) strategy for the existing framework AL4PDE. 
AL4PDE is an active learning framework in which an ensemble of neural networks approximates the solution trajectories of a PDE. The trainings data consits of a specific PDE with the associated parameters and initial conditions, paired with the belonging solution trajectory. The solution is provided by a numerical scheme, depending on the PDE. 
Our proposed method retains the core principle of using model uncertainty - derived from variance across the ensemble - as a measure of sample informativeness.

Previously, the existing AL strategy uses a pool-based sampling strategy, where a large set of samples is derived from the input space and the uncertainty is evaluated. However, this approach requires calculating the uncertainty for each sample in the pool, which is computationally expensive and limits the sampling to the given pool. 

To address this, we proposed a query synthesis approach that samples directly from the input space, eliminating the need for pool creation and allowing for more flexible exploration.
Crucially, in our case, the gradient of uncertainty is available, allowing the method to steer sampling towards globally more informative regions. 

For query synthesis, we use Markov Chain Monte Carlo (MCMC) methods, where the target distribution is the uncertainty.
Our initial strategy used the Metropolis-adjusted Langevin algorithm, which combines both a discretised version of Langevin diffusion and a Metropolis-Hastings step. Thus, it satisfies both the convergence to uncertainty and the gradient information. The chain was used as input to generate training data. However, the results showed that repeated evaluation of rejected samples led to inefficient learning, as the ensemble had already learned from these instances, resulting in limited loss reduction.

To mitigate this, we removed the Metropolis-Hastings step and accepted all proposed states unconditionally. While this improved performance over MALA, it was still inferior to the original pool-based method. To further improve efficiency, we introduced a parallelized approach where multiple chains were initialized with independent proposals and only the final state of each chain was used for training. The chain length was controlled by a hyperparameter.

This parallelized strategy produced results comparable to the original method. However, due to the computational cost of maintaining multiple chains - especially for longer chain lengths - our method became slower than pool-based sampling with large pool sizes. We showed that reducing the chain length per sample mitigated this cost, albeit at the cost of increased loss.

Based on the foundation laid in this work, future improvements should include sampling techniques utilizing an adaptation of MCMC while not being dependent on the Metropolis-Hastings acceptance step. One potential way of achieving this with the current implementation is to start MALA on already trained samples. Therefore, by proposing the next step, the probability of acceptance should be higher, since unknown samples typically have a higher uncertainty than samples already trainend on. However, the problem with duplicates will probably still exist. One further usage could be to not use the chain as input. Instead, use it to create the pool in pool-based methods. This could lead to a pool with samples containing higher uncertainties, makeing future sampling more effective.

More broadly, advanced sampling methods, such as Hamiltonian Monte Carlo (HMC) and Population Monte Carlo (PMC) offer promising directions. 

HMC uses Hamiltonian dynamics to propose transitions in the input space, using both gradients and momentum variables to generate proposals that follow the target distribtution more efficiently. HMC avoids random walk behaviour, allowing faster convergence - particular in high dimensional spaces like the PDE input space. HMC is also the more general form of MALA - or the other way around, MALA is a specific implementation of HMC. However, since HMC still has an acceptance step it could potentially run into the same errors as MALA.

If we want to completly get rid of the Metropolis-Hastings step, PMC would be an alternative. PMC offers a population-based approach where multiple samples are evolved simultaneously, with importance weights guiding resampling and adaptation over iterations. PMC also naturally supports parallelism and could scale well with our existing parallel chain architecture.

In summary, our work presents a viable alternative to pool-based sampling in AL4PDE, offering a gradient-informed and more flexible sampling methods. While still trade-offs in runtime and accuracy remain, our approach lays the groundwork for future improvements in efficient, uncertainty-guided query synthesis for PDE-based learning frameworks.